% Shape: mathematical research article on the AMS article class.
% Evidence for each part: ../evidence/mathematical-writing/ (keys in ../evidence/bibliography.md).
% Compiles with pdflatex. Packages: amsart, amsmath, amssymb, amsthm only.
\documentclass{amsart}
\usepackage{amsmath,amssymb}
\usepackage{amsthm} % amsart already contains amsthm; the line documents the dependency

% Theorem environments in the three AMS styles (AMS 2017).
\theoremstyle{plain}
\newtheorem{theorem}{Theorem}[section]
\newtheorem{proposition}[theorem]{Proposition}
\newtheorem{lemma}[theorem]{Lemma}
\newtheorem{corollary}[theorem]{Corollary}
\theoremstyle{definition}
\newtheorem{definition}[theorem]{Definition}
\newtheorem{example}[theorem]{Example}
\theoremstyle{remark}
\newtheorem{remark}[theorem]{Remark}

% One line naming the results a proof depends on (blueprint practice, Massot 2026).
\newcommand{\uses}[1]{\par\noindent{\normalfont\small\textsc{Uses:} #1.}\par\smallskip}

\begin{document}

\title[Short running title]{A title that names the result}
\author{Author Name}
\address{Department, Institution, City, Country}
\email{author@example.org}

% Put the abstract before \maketitle so amsart sets it in its proper place (AMS 2017).
\begin{abstract}
One to three sentences: what is proved, under which hypotheses, and why it matters.
No citations and no macros of your own; suggested ceiling of about 150 words for a short paper, 300 for a long one.
On arXiv the abstract must stay under 1920 characters.
\end{abstract}

% At least one primary code is required; replace 00-XX by full codes from msc2020.org.
\subjclass[2020]{Primary 00-XX; Secondary 00-XX}
\keywords{first keyword, second keyword}

\maketitle

\section{Introduction}\label{sec:intro}

% Paragraph 1: the result in plain words, before any technical definition (Krantz 2017).
State the main result in one or two plain sentences that a reader outside the
subfield can follow.

\subsection{Main results}\label{subsec:main}
% The theorem comes first and stands alone: one sentence, all hypotheses inside it
% (Halmos 1970, Knuth 1989). If it is too technical here, state an informal version
% labelled as informal and point to the precise one (Tao 2026a).
\begin{theorem}\label{thm:main}
Let $X$ be an object of the kind studied here. If $X$ satisfies condition $(\ast)$,
then $X$ has property $P$.
\end{theorem}

\begin{corollary}\label{cor:main}
A consequence of Theorem~\ref{thm:main} that a reader in a neighbouring field would use.
\end{corollary}

\subsection{Relation to previous work}\label{subsec:previous}
% What was known, what is new, and why no hypothesis can be dropped (Tao 2026a).
Say what was known before, in one sentence per result, with citations~\cite{ref:first}.
Say what this paper adds. Give an example showing that condition $(\ast)$ cannot be removed.

\subsection{Idea of the proof}\label{subsec:idea}
% The milestones and how they combine: the reader's map of the proof (Krantz 2017, Tao 2026b).
The proof has three steps. Proposition~\ref{prop:reduction} reduces Theorem~\ref{thm:main}
to a statement about a simpler object. Lemma~\ref{lem:first} proves that statement for a
special case, and Lemma~\ref{lem:second} extends it. The new difficulty sits in
Lemma~\ref{lem:second}.

\subsection{Organisation}\label{subsec:organisation}
Section~\ref{sec:prelim} fixes notation and recalls known results.
Section~\ref{sec:results} states the milestones and derives Theorem~\ref{thm:main} from them.
Section~\ref{sec:proofs} proves the milestones.

\section{Preliminaries and notation}\label{sec:prelim}
% Notation is fixed once, here, so every later statement can stand alone (Kuteeva 2015).
Throughout, $X$ denotes an object of the kind studied here and $n$ a positive integer.

\begin{definition}\label{def:object}
An object $X$ satisfies \emph{condition $(\ast)$} if it has the defining property.
\end{definition}

\begin{theorem}[Known result~\cite{ref:second}]\label{thm:known}
Quote the known theorem exactly as it is used, with its source.
\end{theorem}

\section{Results}\label{sec:results}
% Milestones first, proofs later: the reader sees the whole argument before the details
% (Krantz 2017, Tao 2026b). Each statement says what it uses.

\begin{proposition}\label{prop:reduction}
If Lemmas~\ref{lem:first} and~\ref{lem:second} hold, then every $X$ satisfying
condition $(\ast)$ has property $P$.
\end{proposition}

\begin{lemma}\label{lem:first}
The special case of property $P$ holds.
\end{lemma}

\begin{lemma}\label{lem:second}
If the special case holds, then the general case holds.
\end{lemma}

\begin{proof}[Proof of Theorem~\ref{thm:main}]
\uses{Proposition~\ref{prop:reduction}, Lemmas~\ref{lem:first} and~\ref{lem:second}}
Combine the three statements: Lemma~\ref{lem:first} gives the special case,
Lemma~\ref{lem:second} extends it, and Proposition~\ref{prop:reduction} concludes.
\end{proof}

\section{Proofs}\label{sec:proofs}
% Each proof opens with a sketch, then numbered steps; see proof.tex and proof.md for the full shape.
% A short proof with a single idea stays a paragraph.

\begin{proof}[Proof of Proposition~\ref{prop:reduction}]
\uses{Definition~\ref{def:object}, Theorem~\ref{thm:known}}
A short argument stays a paragraph: say what is shown, then why, naming each
result used at the step where it is used.
\end{proof}

\begin{proof}[Proof of Lemma~\ref{lem:first}]
\uses{Definition~\ref{def:object}, Theorem~\ref{thm:known}}
\emph{Sketch.} One sentence that gives the idea, so a reader can stop here.
\begin{enumerate}
\item[$\langle 1\rangle 1.$] First claim. \emph{Proof:} by Definition~\ref{def:object}.
\item[$\langle 1\rangle 2.$] Second claim. \emph{Proof:} by $\langle 1\rangle 1$ and Theorem~\ref{thm:known}.
\item[$\langle 1\rangle 3.$] \textsc{Q.e.d.} \emph{Proof:} by $\langle 1\rangle 1$ and $\langle 1\rangle 2$.
\end{enumerate}
\end{proof}

\begin{proof}[Proof of Lemma~\ref{lem:second}]
\uses{Lemma~\ref{lem:first}}
\emph{Sketch.} One sentence that names the new difficulty and how it is overcome.
\begin{enumerate}
\item[$\langle 1\rangle 1.$] Reduction to the special case. \emph{Proof:} by Lemma~\ref{lem:first}.
\item[$\langle 1\rangle 2.$] \textsc{Q.e.d.} \emph{Proof:} by $\langle 1\rangle 1$.
\end{enumerate}
\end{proof}

% Optional ending: open questions or remarks; the field has no fixed ending (Kuteeva 2015).
\section*{Acknowledgements}
Funding and thanks.

\begin{thebibliography}{9}
\bibitem{ref:first} A. Author, \emph{Title of the earlier result}, Journal \textbf{1} (2000), 1-10.
\bibitem{ref:second} B. Author, \emph{Title of the known theorem}, Publisher, City, 2010.
\end{thebibliography}

\end{document}
